Um zu zeigen, dass die orthogonale Gruppe $\mathrm{O}(n)$ kompakt ist, müssen wir zwei Eigenschaften nachweisen: $\mathrm{O}(n)$ ist abgeschlossen und beschränkt in $\mathbb{R}^{n \times n}$, dem Raum der $n \times n$-Matrizen mit den reellen Zahlen als Einträgen.

Zunächst zeigen wir, dass $\mathrm{O}(n)$ abgeschlossen ist. Die Gruppe $\mathrm{O}(n)$ besteht aus allen $n \times n$-Matrizen $A$, die die Bedingung $A^T A = I_n$ erfüllen, wobei $A^T$ die Transponierte von $A$ und $I_n$ die $n \times n$-Einheitsmatrix ist. Die Abbildung $A \mapsto A^T A$ ist stetig, da sie sich aus der Multiplikation und Transposition zusammensetzt, die beide stetige Operationen sind. Die Menge der Matrizen, die die Gleichung $A^T A = I_n$ erfüllen, ist das Urbild der stetigen Abbildung $A \mapsto A^T A$ unter der konstanten Abbildung $A \mapsto I_n$. Da das Urbild eines Punktes unter einer stetigen Abbildung abgeschlossen ist und die Einheitsmatrix $I_n$ ein fester Punkt ist, folgt, dass $\mathrm{O}(n)$ abgeschlossen ist.

Nun zeigen wir, dass $\mathrm{O}(n)$ beschränkt ist. Jede Matrix $A \in \mathrm{O}(n)$ erfüllt $A^T A = I_n$, was bedeutet, dass die Spalten von $A$ orthonormale Vektoren sind. Insbesondere hat jede Spalte von $A$ die Norm 1. Daher sind die Einträge von $A$ durch 1 beschränkt, was bedeutet, dass $\mathrm{O}(n)$ als Teilmenge von $\mathbb{R}^{n \times n}$ beschränkt ist.

Da $\mathrm{O}(n)$ sowohl abgeschlossen als auch beschränkt in $\mathbb{R}^{n \times n}$ ist und $\mathbb{R}^{n \times n}$ ein endlich-dimensionaler Raum ist, folgt aus dem Satz von Heine-Borel, dass $\mathrm{O}(n)$ kompakt ist.